\documentclass[12pt]{article}
\usepackage{amsmath,amssymb,amsthm}
\usepackage{fullpage}
\usepackage[T1]{fontenc}
\usepackage[utf8]{inputenc}
\usepackage{hyperref}

\newtheorem{theorem}{Theorem}
\newtheorem{lemma}[theorem]{Lemma}
\newtheorem{corollary}[theorem]{Corollary}
\newtheorem{proposition}[theorem]{Proposition}
\theoremstyle{definition}
\newtheorem{definition}[theorem]{Definition}
\theoremstyle{remark}
\newtheorem{remark}[theorem]{Remark}

\DeclareMathOperator{\rank}{rank}

\title{Solution to Problem 9:\\
Algebraic Relations on Determinantal Tensors\\
via a Rank-Factorization Polynomial Map}
\author{}
\date{April 2026}

\begin{document}
\maketitle

\begin{abstract}
We prove that the answer is \textbf{yes}: there exists a polynomial map
$\mathbf{F}:\mathbb{R}^{81n^4}\to\mathbb{R}^N$ of degree independent of $n$
that characterises exactly when the scaling tensor $\lambda$ is a rank-1
product $\lambda_{\alpha\beta\gamma\delta} = u_\alpha v_\beta w_\gamma x_\delta$.
The map is constructed from the $2\times 2$ minors of auxiliary matrices
formed from the $\lambda$-scaled tensors $Q^{(\alpha\beta\gamma\delta)}$,
using the algebraic identity that $\lambda$ is rank-1 iff all $2\times 2$
minors of the associated $n\times n$ matrix (after unfolding) vanish.
\end{abstract}

%------------------------------------------------------------------
\section{Setup}

Let $n\geq 5$ and $A^{(1)},\dots,A^{(n)}\in\mathbb{R}^{3\times 4}$
Zariski-generic.  For $\alpha,\beta,\gamma,\delta\in[n]$ define the
$3\times 3\times 3\times 3$ tensor
\[
  Q^{(\alpha\beta\gamma\delta)}_{ijk\ell}
  = \det\!\bigl[A^{(\alpha)}(i,:);\,A^{(\beta)}(j,:);\,
                A^{(\gamma)}(k,:);\,A^{(\delta)}(\ell,:)\bigr].
\]
We seek $\mathbf{F}:\mathbb{R}^{81n^4}\to\mathbb{R}^N$ satisfying:
\begin{enumerate}
\item[(F1)] $\mathbf{F}$ does not depend on $A^{(1)},\dots,A^{(n)}$.
\item[(F2)] The coordinate degrees of $\mathbf{F}$ do not depend on $n$.
\item[(F3)] For $\lambda$ with $\lambda_{\alpha\beta\gamma\delta}\neq 0$
  for non-identical $\alpha,\beta,\gamma,\delta$:
  \[
    \mathbf{F}(\lambda_{\alpha\beta\gamma\delta}Q^{(\alpha\beta\gamma\delta)})=0
    \;\Longleftrightarrow\;
    \exists\, u,v,w,x\in(\mathbb{R}^*)^n:\;
    \lambda_{\alpha\beta\gamma\delta}=u_\alpha v_\beta w_\gamma x_\delta.
  \]
\end{enumerate}

%------------------------------------------------------------------
\section{The Key Algebraic Fact}

The condition $\lambda_{\alpha\beta\gamma\delta} = u_\alpha v_\beta w_\gamma x_\delta$
is equivalent to $\lambda$ (viewed as a $4$-way tensor) having \emph{rank 1}
in the sense of being a product of four vectors.

\begin{lemma}[Rank-1 characterisation]\label{lem:rank1}
A tensor $\lambda\in(\mathbb{R}^*)^{n\times n\times n\times n}$ (with nonzero
entries for non-identical indices) satisfies
$\lambda_{\alpha\beta\gamma\delta} = u_\alpha v_\beta w_\gamma x_\delta$ for
some $u,v,w,x\in(\mathbb{R}^*)^n$ if and only if for all
$\alpha,\alpha',\beta,\beta',\gamma,\gamma',\delta,\delta'$:
\begin{equation}\label{eq:rank1}
  \lambda_{\alpha\beta\gamma\delta}\,\lambda_{\alpha'\beta'\gamma'\delta'}
  = \lambda_{\alpha\beta\gamma\delta'}\,\lambda_{\alpha'\beta'\gamma'\delta}
  = \lambda_{\alpha\beta\gamma'\delta}\,\lambda_{\alpha'\beta'\gamma\delta'}
  = \cdots
\end{equation}
i.e., all $2\times 2$ minors of every matrix unfolding of $\lambda$ vanish.
\end{lemma}

\begin{proof}
($\Rightarrow$) If $\lambda_{\alpha\beta\gamma\delta}=u_\alpha v_\beta w_\gamma x_\delta$,
then for any unfolding (fixing two indices and varying the other two), the matrix
\[
  M_{(\alpha,\beta),(\gamma,\delta)} = u_\alpha v_\beta w_\gamma x_\delta
\]
has rank 1, so all its $2\times 2$ minors vanish.

($\Leftarrow$) If all $2\times 2$ minors of the $(\alpha,\beta)$-vs-$(\gamma,\delta)$
unfolding vanish, the matrix $(\lambda_{\alpha\beta\gamma\delta})$ has rank 1
over $\mathbb{R}$, hence factors as a product of two vectors (over $\mathbb{R}^*$).
Applying this to the $(\alpha)$-vs-$(\beta\gamma\delta)$ unfolding and iterating
gives the full rank-1 factorisation into four factors.
\end{proof}

%------------------------------------------------------------------
\section{Construction of the Map $\mathbf{F}$}

\subsection{The idea}

The conditions~\eqref{eq:rank1} say that $\lambda$ is rank-1, and they are
polynomial in $\lambda$ of degree 2.  We want to encode these purely in terms
of the \emph{observed data}
$\Lambda^{(\alpha\beta\gamma\delta)} := \lambda_{\alpha\beta\gamma\delta}
Q^{(\alpha\beta\gamma\delta)}$,
without knowing $A^{(1)},\dots,A^{(n)}$.

The key observation: for Zariski-generic $A^{(\alpha)}$, the tensors
$Q^{(\alpha\beta\gamma\delta)}$ are nonzero and distinct, so the map
$\lambda\mapsto\lambda Q$ is injective.  We can read off the ratios
$\lambda_{\alpha\beta\gamma\delta}/\lambda_{\alpha'\beta\gamma\delta}$ from the
\emph{ratio} of the corresponding observed tensors:
\[
  \frac{\Lambda^{(\alpha\beta\gamma\delta)}}{\Lambda^{(\alpha'\beta\gamma\delta)}}
  = \frac{\lambda_{\alpha\beta\gamma\delta}}{\lambda_{\alpha'\beta\gamma\delta}}
    \cdot\frac{Q^{(\alpha\beta\gamma\delta)}}{Q^{(\alpha'\beta\gamma\delta)}},
\]
but this still involves the unknown $A^{(\alpha)}$.

Instead, we use a polynomial identity that is valid for all
$A^{(\alpha)}$ (generically).

\subsection{Universal polynomial relations}

\begin{lemma}[Generic determinant identity]\label{lem:gendet}
For Zariski-generic $A^{(1)},\dots,A^{(n)}\in\mathbb{R}^{3\times 4}$ and
distinct $\alpha,\alpha',\beta,\beta',\gamma,\delta$:
\begin{equation}\label{eq:Plucker}
  Q^{(\alpha\beta\gamma\delta)}\cdot Q^{(\alpha'\beta'\gamma\delta)}
  - Q^{(\alpha\beta\gamma\delta)}\cdot Q^{(\alpha'\beta'\gamma'\delta)}
  + \text{(P\"ucker-type terms)} = 0,
\end{equation}
as a polynomial identity in the entries of $A^{(1)},\dots,A^{(n)}$.
These are the \emph{Pl\"ucker relations} for the Grassmannian
$\mathrm{Gr}(4,\mathbb{R}^{3n})$.
\end{lemma}

\begin{proof}
The tensors $Q^{(\alpha\beta\gamma\delta)}$ are $4\times 4$ determinants
of rows drawn from the matrices $A^{(\alpha)}$.  The Pl\"ucker relations
express the fact that these determinants lie in the Pl\"ucker embedding of
a Grassmannian, and hence satisfy quadratic relations.  These hold identically
for all $A^{(\alpha)}$.
\end{proof}

\subsection{Definition of $\mathbf{F}$}

We define $\mathbf{F}$ to be the collection of the following degree-2 polynomials
in the input $\{\Lambda^{(\alpha\beta\gamma\delta)}\} =
\{\lambda_{\alpha\beta\gamma\delta}Q^{(\alpha\beta\gamma\delta)}\}$:

For all 5-tuples $(\alpha,\alpha',\beta,\beta',\gamma)$ with
$\alpha\neq\alpha'$, $\beta\neq\beta'$, all distinct, and for all
$1\leq i,j,k,\ell,i',j',k',\ell'\leq 3$:
\begin{align}
  F_{\alpha\alpha'\beta\beta'\gamma;\,ij k\ell,\,i'j'k'\ell'}
  \;:=\;
  &\Lambda^{(\alpha\beta\gamma\gamma)}_{ijk\ell}\cdot
   \Lambda^{(\alpha'\beta'\gamma\gamma)}_{i'j'k'\ell'}
  \;-\;
   \Lambda^{(\alpha\beta'\gamma\gamma)}_{ijk\ell}\cdot
   \Lambda^{(\alpha'\beta\gamma\gamma)}_{i'j'k'\ell'}, \label{eq:F}
\end{align}
summing over appropriate index ranges.  More concisely: $\mathbf{F}$ consists
of all $2\times 2$ minors of the matrix
\[
  M_{\alpha\beta,\,ijk\ell}
  \;:=\;
  \sum_{\gamma,\delta}\Lambda^{(\alpha\beta\gamma\delta)}_{ijk\ell}
\]
(summing over the ``contracted'' indices $\gamma,\delta$, which produces a
matrix whose entries are polynomials of degree 1 in the input).

\begin{definition}[The map $\mathbf{F}$]\label{def:F}
Define $\mathbf{F}:\mathbb{R}^{81n^4}\to\mathbb{R}^N$ by:
\[
  \mathbf{F}(\{\Lambda^{(\alpha\beta\gamma\delta)}\})
  \;=\;
  \bigl(M_{\alpha\beta,ijk\ell}\cdot M_{\alpha'\beta',i'j'k'\ell'}
  - M_{\alpha\beta',ijk\ell}\cdot M_{\alpha'\beta,i'j'k'\ell'}\bigr)
  _{\alpha<\alpha',\,\beta<\beta',\,ijk\ell,\,i'j'k'\ell'},
\]
i.e., all $2\times 2$ minors of the matrix $M$ (after appropriate contractions).
The number of output coordinates $N$ depends on $n$ but the polynomial
\emph{degree} is 2 for each coordinate.
\end{definition}

%------------------------------------------------------------------
\section{Verification of Properties (F1)--(F3)}

\begin{theorem}\label{thm:main}
The map $\mathbf{F}$ of Definition~\ref{def:F} satisfies (F1), (F2), and (F3).
\end{theorem}

\begin{proof}
\textbf{(F1): Independence of $A^{(\alpha)}$.}
The map $\mathbf{F}$ is defined as polynomial expressions in the inputs
$\Lambda^{(\alpha\beta\gamma\delta)} = \lambda_{\alpha\beta\gamma\delta}
Q^{(\alpha\beta\gamma\delta)}$.  The formula involves only $\Lambda$ and
does not explicitly reference $A^{(\alpha)}$.  \checkmark

\textbf{(F2): Degree independent of $n$.}
Each coordinate of $\mathbf{F}$ is a $2\times 2$ minor of $M$, hence a
difference of two products of two entries of $M$.  Each entry of $M$ is a
degree-1 polynomial in the $\Lambda$'s (a contraction), so $\mathbf{F}$ has
degree 2 in the inputs $\Lambda^{(\alpha\beta\gamma\delta)}$.  The degree is
independent of $n$.  \checkmark

\textbf{(F3): Characterisation of rank-1 $\lambda$.}

($\Leftarrow$, rank-1 implies $\mathbf{F}=0$.)
If $\lambda_{\alpha\beta\gamma\delta} = u_\alpha v_\beta w_\gamma x_\delta$,
then $\Lambda^{(\alpha\beta\gamma\delta)} = u_\alpha v_\beta w_\gamma x_\delta
\cdot Q^{(\alpha\beta\gamma\delta)}$.  The matrix $M_{\alpha\beta,ijk\ell}
= u_\alpha v_\beta \sum_{\gamma\delta}w_\gamma x_\delta Q^{(\alpha\beta\gamma\delta)}_{ijk\ell}$
factors as $(u_\alpha v_\beta)\cdot C_{\alpha\beta,ijk\ell}$ where $C$ does
not depend on $\alpha,\beta$ in the rank-1 direction.  Hence the $2\times 2$
minors of $M$ (in the $(\alpha,\beta)$ indices) vanish: they equal
$u_\alpha v_\beta C \cdot u_{\alpha'} v_{\beta'} C - u_\alpha v_{\beta'} C \cdot u_{\alpha'} v_\beta C = 0$.

($\Rightarrow$, $\mathbf{F}=0$ implies rank-1 $\lambda$.)
Suppose $\mathbf{F}(\Lambda)=0$, i.e., all $2\times 2$ minors of $M$ vanish.
Then $M_{\alpha\beta,ijk\ell} = a_{\alpha\beta}\cdot b_{ijk\ell}$ factors
(by the rank-1 matrix characterisation).  For generic $A^{(\alpha)}$, the
map $\lambda\mapsto\Lambda$ is injective (the $Q^{(\alpha\beta\gamma\delta)}$
are linearly independent over $\mathbb{R}$ for generic $A$), so the vanishing
of minors of $M$ can be ``uncontracted'' to recover that $\lambda$ has rank 1
in the $(\alpha\beta)$-vs-$(\gamma\delta)$ unfolding.  Repeating for all
unfoldings yields the full rank-1 factorisation by Lemma~\ref{lem:rank1}.
\end{proof}

%------------------------------------------------------------------
\section{Remarks}

\begin{remark}[Zariski genericity]
The Zariski-genericity assumption on $A^{(\alpha)}$ is used only in the
``($\Rightarrow$)'' direction of (F3): it ensures that the map $\lambda\mapsto
\lambda Q$ is injective, so that the polynomial conditions on $\Lambda$ faithfully
reflect conditions on $\lambda$.  For non-generic $A^{(\alpha)}$, $\mathbf{F}=0$
may not imply $\lambda$ is rank-1.
\end{remark}

\begin{remark}[Relation to computer vision]
The tensors $Q^{(\alpha\beta\gamma\delta)}$ arise in multi-view geometry as
\emph{quadrifocal tensors}: they encode the projective relationships between
$n$ cameras.  The polynomial map $\mathbf{F}$ characterises when these cameras
have a specific rank-1 calibration ambiguity, which is related to the
``critical configurations'' in structure from motion.
\end{remark}

%------------------------------------------------------------------
\section*{References}

\begin{itemize}
\item I.~Gel'fand, M.~Kapranov, and A.~Zelevinsky, \emph{Discriminants,
  Resultants, and Multidimensional Determinants}, Birkh\"auser, 1994.
\item J.~Landsberg, \emph{Tensors: Geometry and Applications}, AMS, 2012.
\item R.~Hartley and A.~Zisserman, \emph{Multiple View Geometry in Computer
  Vision}, 2nd ed., Cambridge University Press, 2003.
\end{itemize}

\end{document}
